% Companion to CM_j0_degree10.tex. No odd-degree parity argument is used.
\begin{lemma}\label{lem:decic242-order3-stabilizers}
Put $X=X_1(2,42)$. Let $\overline N$ be a noncuspidal effective
divisor of degree $9$ on $X_{\mathbb F_5}$, fixed by a rational
level automorphism $h$ of order $3$.
There is no effective divisor $D$ on $X$, defined over $\mathbb Q$,
of degree $10$, with entirely noncuspidal support and reduction
$\overline N+\overline c$, where $\overline c$ is a cusp.
\end{lemma}
\begin{proof}
The rational Jacobian of $X$ has rank zero by
\cite[Theorem~3.1]{DEHMZ}.
Its projective index is $1152$, so the bound in
\cite[(2.1)]{NV} gives
\[
 \operatorname{gon}_{\overline{\mathbb Q}}X
 \geq\frac{325\cdot1152}{32768}>11.
\]
Suppose that $D$ exists. Lift $\overline c$ to a cusp $c$ over
an extension unramified at $5$.
The class
\[
 [hD-D-hc+c]
\]
is torsion by rank zero and Manin--Drinfeld. Its reduction is
zero, since $h\overline N=\overline N$.
Injectivity of torsion reduction over an unramified extension of
$\mathbb Q_5$ gives
\[
 hD+c\sim D+hc.
\]
These are effective divisors of degree $11$, below geometric
gonality. Hence they are equal.
Comparing their noncuspidal parts gives $hD=D$.

A geometric fixed point of $h$ would give an elliptic
automorphism of order $3$ or $6$ and would therefore have
$j$-invariant $0$.
Every closed point in the support of $D$ has degree at most $10$.
Lemma~\ref{lem:decic242-no-j0} excludes such a point with
$j$-invariant $0$. Thus $h$ acts freely on the geometric support
of $D$. Its orbits have length $3$, and $hD=D$ makes the
multiplicities constant on each orbit. It follows that $3$
divides $\deg D=10$, a contradiction.
\end{proof}

\noindent\emph{Application to the three degree-$9$ supports.}
Let
\[
 k=\mathbb F_5[\alpha]/(\alpha^9+2\alpha^3+\alpha+3).
\]
For $B=\sum_{i=0}^8B_i\alpha^i$ and
$C=\sum_{i=0}^8C_i\alpha^i$, use the level structure
\[
 E:\ y^2=x^3+Cx^2+B(1-B-C)x,
 \qquad P=(0,0),\quad Q=(B,B).
\]
The three remaining degree-$9$ supports have coefficient vectors
\[
\begin{array}{c|c|c}
\text{case}&(B_i)&(C_i)\\ \hline
1777&(0,0,0,0,2,0,2,1,4)&(4,4,3,1,0,0,4,2,2)\\
1780&(0,0,1,0,1,3,4,1,2)&(3,4,1,0,4,1,1,4,4)\\
1783&(0,0,0,1,0,1,0,2,3)&(4,2,3,2,3,3,0,2,1).
\end{array}
\]
Put
\[
 h_1(P,Q)=(21Q,P+25Q),\qquad
 h_2(P,Q)=(P+21Q,P+4Q).
\]
Both have order $3$. On case $1777$, $h_1$ agrees with the
third power of Frobenius. On cases $1780$ and $1783$, $h_2$
agrees with the third and sixth powers of Frobenius,
respectively. Each closed degree-$9$ divisor is therefore fixed
by the indicated automorphism, and
Lemma~\ref{lem:decic242-order3-stabilizers} excludes all three.

The complete automorphism and cusp-action enumeration is in
\path{next242_cheap_stabilizers.py}; it checks all $36$ level
automorphisms on all $111$ input supports and all $72$ geometric
cusps. It removes exactly these three cases.
The three displayed Frobenius identities and the exact order-$3$
level identities are independently verified in Magma by
\path{next242_order3_phases.m}; the output is
\path{degree9_order3_native_phases.m}.
The remaining $108$ supports, with their original case IDs and
field polynomials, are stored in
\path{after_cheap_stabilizers_unified.m}.
